\newcommand{\GBMEMRate}{0.500}
\newcommand{\GBMEMWeakRate}{0.999}
\newcommand{\GBMMilRate}{0.991}
\newcommand{\GBMMilWeakRate}{1.000}
\newcommand{\GBMAnalyticWeakRate}{0.9993}
\newcommand{\GBMFitWindow}{$N=64,128,256,512,1024$}
\newcommand{\GBMExactMean}{105.1987}
\newcommand{\GBMExactVariance}{1045.0120}
\newcommand{\GBMExactMeanCI}{[105.0403, 105.3571]}
\newcommand{\NARate}{0.491}
\newcommand{\NAFitWindow}{$N=8,16,32,64$}
\newcommand{\NAExactMean}{105.08796}
\newcommand{\NAExactMeanCI}{[104.89962, 105.27629]}
\newcommand{\NAPayoff}{14.53804}
\newcommand{\NAPayoffCI}{[14.40515, 14.67093]}
\newcommand{\NAWeakInterpretation}{在 $N=8,16,32$ 时，区间不含零，因此这些较粗网格的收益差可以从抽样噪声中分辨。在 $N=64,128,256$ 时，区间包含零，现有样本无法可靠判断偏差符号。这不证明偏差为零，也不足以为完整网格范围指定弱阶。}
\newcommand{\NAReferenceInterpretation}{相邻参考网格的平均绝对差从 2048 与 4096 步时的 0.05679，降至 4096 与 8192 步时的 0.04016，后者约为前者的 0.707 倍，与半阶缩放相容。在 $N=8,16,32,64$ 的拟合窗口内，0.04016 约占粗网格到最细参考强差的 3.2\% 至 8.9\%；在 $N=128,256$ 时，这一比例升至 12.6\% 和 18.0\%。因此，更细的粗网格相对更受参考分辨率影响，主要斜率使用参考影响相对较小的四个网格。这些比较衡量参考加密的敏感性，不构成相对于未知精确解的误差上界。}
